\documentclass[noauthor]{ximera}

\title{Week 2}
\author{Math 1193 Design Team}

\begin{document}
\begin{abstract}
    Solutions for Week 2.
\end{abstract}
\maketitle

\begin{problem}
Say you sell shoes at \$120 per pair.

\begin{enumerate}
    \item How much do you 
    make if you sell 0 pairs of shoes? 1? 2? 3? 
    \item How much do you 
    make if you sell 32 pairs of shoes?
    \item How much do you make if you sell \(x\)
    pairs of shoes?
    \item Write the formula for a function \(R\)
    which takes in a number of pairs of shoes and
    returns the revenue made.
    \item Graph your function \(R\).
    \item What does the slope of \(R\) represent
    in the problem?
\end{enumerate}

\begin{explanation}

    \begin{table}[h]
\begin{tabular}{|c|c|} \hline
Pairs of shoes sold & Revenue (\$) \\ \hline
0 & 0 \\
1 & 120 \\
2 & 240 \\
3 & 360 \\
32 & 3840 \\
$x$ & $120x$ \\ \hline
\end{tabular}
\end{table}
\begin{itemize}
    \item \(R(x) = 120x\)
    
    Notice that when writing a function, we use function notation to indicate the independent variable (the input).
    Another possible answer could be \(R(p) = 120p\). The important thing is that your formula should say what the input is and how to get from the input to the output.

    \item 
    \begin{center}
    \includegraphics[width=10cm]{revenueGraph.png}
    \end{center}

    This graph is very steep.

    \item The slope represents the revenue per pair of shoes sold. 
    
    Since slope is a change in \(y\) over a change in \(x\), it should be a change in outputs (revenue) over a change in inputs (pairs of shoes sold).
\end{itemize}
\end{explanation}

\end{problem}

\begin{problem}
Say it costs you \$30 in production costs to make a pair
of shoes and \$2000 per month to rent out the shoe-making
facility.

\begin{enumerate}
    \item How much does it cost to 
    make 0 pairs of shoes per month? 1? 2? 3? 
    \item How much does it cost to 
    make 12 pairs of shoes per month?
    \item How much does it cost to 
    make \(x\) pairs of shoes per month? 
    \item Write the formula for a cost function \(C\)
    which takes in a number of pairs of shoes and
    returns the cost per month to make them.
    \item Graph your function \(C\).
    \item What does the slope of \(C\) represent
    in the problem?
\end{enumerate}

\begin{explanation}

\begin{table}[h]
\begin{tabular}{|c|c|}
    \hline
Pairs of shoes made & Cost (\$) \\ \hline
0 & 2000 \\
1 & 2030 \\
2 & 2060 \\
3 & 2090 \\
12 & 2360 \\
$x$ & $30x + 2000$ \\ \hline
\end{tabular}
\end{table}
\begin{itemize}
    \item \(C(x) = 30x + 2000\)
    
    Notice that when writing a function, we use function notation to indicate the independent variable (the input).
    Another possible answer could be \(C(n) = 30n + 2000\). The important thing is that your formula should say what the input is and how to get from the input to the output.

    \item 
    \begin{center}
    \includegraphics[width=10cm]{costGraph.png}
    \end{center}

    This graph is very steep.

    \item The slope represents the cost per pair of shoes sold. 
    
    Since slope is a change in \(y\) over a change in \(x\), it should be a change in outputs (cost) over a change in inputs (pairs of shoes sold).

    If you already know about the formulas of lines, you can note that the slope is 30, and has nothing to do with the 2000. Since slope only measures a rate of change,
    and the 2000 is a fixed cost, this makes sense.
\end{itemize}
\end{explanation}


\end{problem}

\begin{problem}
Consider the function \(f(x) = \frac{1}{2}x + 5\).

\begin{enumerate}
    \item Write down a word problem about some
    quantity that can be represented by \(f(x)\).
    Be sure to define the inputs and outputs.
    \item What are the \(x\)- and \(y\)-intercepts
    of the function? What do those represent
    in your problem?
    \item What is the slope of \(f\)? What does 
    that represent in your problem?
\end{enumerate}

\begin{explanation}
Here's one possibility:

\begin{enumerate}
    \item You and your kid brother are running a lemonade stand. You start the day with \$5 in cash on hand and sell each glass of lemonade for \$0.50.
    
    Write down a function \(f\) whose input is the number of glasses sold and whose output is the amount of cash you have on hand.

    \item The \(x\)-intercept is \((-10, 0)\). In this context, that means that if you sell \(-10\) glasses of lemonade, you'll have \$0 in cash on hand. Not a super relevant scenario.
    
    The \(y\)-intercept is \((0, 5)\). In this context, if you sell 0 glasses of lemonade, you have \$5 on hand. You could also think about this as the amount of cash you start with.

    \item The slope of \(f\) is \(\frac{1}{2}\). In this case, the slope is the revenue per glass of lemonade sold.
\end{enumerate}
\end{explanation}
\end{problem}

\begin{problem}
Say \(f(x) = 2x - g(x)\) and \(g(x) = 5 + x\).

\begin{enumerate}
    \item What is the slope of \(f\)?
    \item What is the \(y\)-intercept?
    \item What is \(f(2)\)?
    \item Graph \(f\).
\end{enumerate}

\begin{explanation}
    First, it's helpful to find a formula for \(f(x)\) that does not involve \(g(x)\). 
    \begin{align*}
        f(x) & = 2x - g(x) \\
        & = 2x - (5 + x) \\
        & = 2x - 5 - x \\
        & = x - 5
    \end{align*}
    Notice that when substituting the expression for \(g(x)\), we used parentheses. Often, those parentheses don't matter when substituting, but 
    in this case, they do! This is because of the negative sign, which must be distributed to both the 5 and the \(x\). 
\begin{enumerate}
    \item The slope of \(f\) is 1. If you know that the coefficient of the \(x\) term in a linear function is the slope, that's great!
    
    If not, you could always do 
    \begin{align*}
        \frac{\text{change in }y}{\text{change in }x} & = \frac{y_2 - y_1}{x_2 - x_1} \\
        & = \frac{f(x_2) = f(x_1)}{x_2 - x_1} \\
        & = \frac{f(1) - f(0)}{1 - 0} \\
        & = \frac{-4 - (-5)}{1} \\
        & = \frac{1}{1} \\
        & = 1
    \end{align*}

    We chose \(x_2 = 1\) and \(x_1 = 0\) just so the denominator works out nicely.
    
    \item Again, if you know that the \(y\)-intercept of a linear function corresponds to the constant term, that's great! You can then immediately tell it's \((0, -5)\)
    
    If not, remember that the \(y\)-intercept of any function is \((0, f(0))\), that is, you plug in 0 for \(x\). Since \(f(0) = -5\), we get \((0, -5)\).

    \item \(f(2) = 2 - 5 = 3\). Remember that the number inside the parentheses in \(f(2)\) represents the \(x\)-value we're plugging into our formula.
    
    \item \begin{center}
    \includegraphics[width=10cm]{subGraph.png}
    \end{center}

    Again, you could use what you know about linear functions to graph this. You know the \(y\)-intercept is \((0, -5)\) and you know the slope is 1, so you can go up 1 and over 1
    and connect your dots. 

    However, you can also do this just by plugging in test \(x\)-values and seeing what you get. 

\end{enumerate}
\end{explanation}
\end{problem}

\end{document}
